\documentclass[11pt]{article}
\usepackage[margin=1in]{geometry}
\usepackage{amsmath,amssymb}
\setlength{\parindent}{0pt}
\setlength{\parskip}{5pt}
\newcommand{\norm}[1]{\left\lVert#1\right\rVert}
\title{Assignment 2}
\author{00103335: Deep Learning and Reinforcement Learning}
\date{Fall 2026}
\begin{document}
\maketitle
All norms are Euclidean. No assumptions beyond those in the questions
are needed; in particular, no second derivatives or Lipschitz gradient
assumptions are used.

\section*{Problem 1}
\textbf{(i)$\Rightarrow$(ii).}
For $0<t<1$, convexity gives
\[
 f(x+t(y-x))=f((1-t)x+ty)\leq(1-t)f(x)+tf(y).
\]
Subtract $f(x)$ and divide by $t>0$:
\[
 \frac{f(x+t(y-x))-f(x)}t\leq f(y)-f(x).
\]
Differentiability gives
\[
 \nabla f(x)^\top(y-x)
 =\lim_{t\downarrow0}\frac{f(x+t(y-x))-f(x)}t
 \leq f(y)-f(x),
\]
which is (ii).

\textbf{(ii)$\Rightarrow$(i).}
Let $z=\lambda x+(1-\lambda)y$. Apply (ii) with base point $z$
and endpoints $x,y$, multiply by $\lambda,1-\lambda$, and add:
\[
 \lambda f(x)+(1-\lambda)f(y)
 \geq f(z)+\nabla f(z)^\top
 [\lambda(x-z)+(1-\lambda)(y-z)]=f(z).
\]

\textbf{(ii)$\Rightarrow$(iii).}
Write (ii) in both directions:
\[
 f(y)-f(x)\geq\nabla f(x)^\top(y-x),\qquad
 f(x)-f(y)\geq\nabla f(y)^\top(x-y).
\]
Adding cancels the function values and gives
\[
 0\geq\nabla f(x)^\top(y-x)+\nabla f(y)^\top(x-y),
\]
which is (iii).

\textbf{(iii)$\Rightarrow$(ii).}
For $x\ne y$, set $h=y-x$ and $g(t)=f(x+th)$.
The mean value theorem gives $g(1)-g(0)=g'(c)$ for some $c\in(0,1)$.
By (iii) at $x,x+ch$,
\[
 (\nabla f(x+ch)-\nabla f(x))^\top(ch)\geq0,
 \qquad g'(c)\geq\nabla f(x)^\top h.
\]
This proves (ii); for $x=y$ it is equality.

\section*{Problem 2}
Define $g(x)=f(x)-\frac\mu2\norm x^2$, so
$\nabla g(x)=\nabla f(x)-\mu x$.

\textbf{(i)$\Leftrightarrow$(ii).}
Set $z=\alpha x+(1-\alpha)y$. The convexity inequality for $g$ is
\[
 f(z)-\frac\mu2\norm z^2
 \leq\alpha f(x)+(1-\alpha)f(y)
 -\frac\mu2\big(\alpha\norm x^2+(1-\alpha)\norm y^2\big).
\]
The quadratic difference expands as
\begin{align*}
 \alpha\norm x^2+(1-\alpha)\norm y^2-\norm z^2
 &=(\alpha-\alpha^2)\norm x^2
   +\big((1-\alpha)-(1-\alpha)^2\big)\norm y^2\\
 &\quad-2\alpha(1-\alpha)x^\top y\\
 &=\alpha(1-\alpha)\big(\norm x^2+\norm y^2-2x^\top y\big)\\
 &=\alpha(1-\alpha)\norm{y-x}^2.
\end{align*}
Moving $\mu\norm z^2/2$ to the right therefore gives
\[
 f(z)\leq\alpha f(x)+(1-\alpha)f(y)
 -\frac\mu2\alpha(1-\alpha)\norm{y-x}^2,
\]
which is (ii). All the algebraic steps are reversible.

\textbf{(i)$\Leftrightarrow$(iii).}
By Problem 1, convexity of $g$ is equivalent to
$g(y)\geq g(x)+\nabla g(x)^\top(y-x)$.
Substituting the definitions gives
\[
 f(y)-\frac\mu2\norm y^2
 \geq f(x)-\frac\mu2\norm x^2
 +(\nabla f(x)-\mu x)^\top(y-x).
\]
Rearranging yields
\begin{align*}
 f(y)&\geq f(x)+\nabla f(x)^\top(y-x)
 +\frac\mu2\big(\norm y^2-\norm x^2-2x^\top(y-x)\big)\\
 &=f(x)+\nabla f(x)^\top(y-x)+\frac\mu2\norm{y-x}^2,
\end{align*}
because the parenthesis equals
$\norm y^2+\norm x^2-2x^\top y=\norm{y-x}^2$.
Reversing these steps proves the converse.

\textbf{(i)$\Leftrightarrow$(iv).}
Again by Problem 1, convexity of $g$ is equivalent to
\[
 0\leq(\nabla g(y)-\nabla g(x))^\top(y-x)
 =(\nabla f(y)-\nabla f(x))^\top(y-x)-\mu\norm{y-x}^2,
\]
which is (iv). Thus all four conditions are equivalent.

\section*{Problem 3}
Use the numbering in the question. Write $h=y-x$ and
$a=\nabla f(y)-\nabla f(x)$.

\textbf{(ii)} Problem 2(iv) and Cauchy--Schwarz give
\[
 \mu\norm h^2\leq a^\top h\leq\norm a\norm h.
\]
For $h\ne0$, divide by $\norm h$ to obtain $\norm a\geq\mu\norm h$;
the result is immediate for $h=0$.

\textbf{(iii)} Apply strong convexity with base point $y$:
\[
 f(x)\geq f(y)+\nabla f(y)^\top(x-y)+\frac\mu2\norm h^2.
\]
Rearranging and completing the square gives
\begin{align*}
 f(y)-f(x)-\nabla f(x)^\top h
 &\leq a^\top h-\frac\mu2\norm h^2\\
 &=\frac1{2\mu}\norm a^2-\frac\mu2\norm{h-a/\mu}^2
 \leq\frac1{2\mu}\norm a^2.
\end{align*}

\textbf{(iv)} By Cauchy--Schwarz and (ii),
\[
 a^\top h\leq\norm a\norm h\leq\frac1\mu\norm a^2.
\]

\section*{Problem 4}
\textbf{(a)} Strong convexity at $x_t$ implies
\[
 f(x^\star)\geq f(x_t)+\nabla f(x_t)^\top(x^\star-x_t)
 +\frac\mu2 d_t^2.
\]
Rearrangement gives
$\nabla f(x_t)^\top(x_t-x^\star)\geq\Delta_t+\frac\mu2d_t^2$.

\textbf{(b)} Expanding the update and using (a) and the stated bound,
\begin{align*}
 d_{t+1}^2
 &=d_t^2-2\eta_t\nabla f(x_t)^\top(x_t-x^\star)
 +\eta_t^2\norm{\nabla f(x_t)}^2\\
 &\leq(1-\mu\eta_t)d_t^2-2\eta_t\Delta_t+\eta_t^2L^2.
\end{align*}

\textbf{(c)} Substitute $\eta_t=2/[\mu(t+1)]$ into (b):
\[
 d_{t+1}^2\leq\frac{t-1}{t+1}d_t^2
 -\frac4{\mu(t+1)}\Delta_t+\frac{4L^2}{\mu^2(t+1)^2}.
\]
Multiplying by $\mu t(t+1)/4$ and rearranging yields
\[
 t\Delta_t\leq\frac\mu4
 \big[t(t-1)d_t^2-t(t+1)d_{t+1}^2\big]
 +\frac{tL^2}{\mu(t+1)}.
\]
To sum, shift the index $s=t+1$ in the negative term:
\begin{align*}
 &\sum_{t=1}^T t(t-1)d_t^2-\sum_{t=1}^T t(t+1)d_{t+1}^2\\
 &\qquad=\sum_{t=1}^T t(t-1)d_t^2
 -\sum_{s=2}^{T+1}s(s-1)d_s^2
 =-T(T+1)d_{T+1}^2.
\end{align*}
The terms with indices $2,\ldots,T$ cancel and the initial coefficient
$1(1-1)$ is zero. Thus
\[
 \sum_{t=1}^T t\Delta_t
 \leq-\frac\mu4 T(T+1)d_{T+1}^2
 +\frac{L^2}\mu\sum_{t=1}^T\frac t{t+1}
 \leq\frac{TL^2}\mu.
\]

\textbf{(d)} The weights $w_t=2t/[T(T+1)]$ are nonnegative and satisfy
\[
 \sum_{t=1}^T w_t=\frac2{T(T+1)}\sum_{t=1}^Tt
 =\frac2{T(T+1)}\frac{T(T+1)}2=1.
\]
Since $\overline x_T=\sum_{t=1}^T w_t x_t$, convexity gives
$f(\overline x_T)\leq\sum_{t=1}^T w_t f(x_t)$.
Subtracting $f(x^\star)=\sum_{t=1}^T w_t f(x^\star)$ and using (c),
\[
 f(\overline x_T)-f(x^\star)
 \leq\sum_{t=1}^T w_t\Delta_t
 =\frac2{T(T+1)}\sum_{t=1}^Tt\Delta_t
 \leq\boxed{\frac{2L^2}{\mu(T+1)}}=O(T^{-1}).
\]
\end{document}
